% Annexe technique du bloc IS (note8p, version 4 : taux de 25 %, CIR supprimé). Chiffres : note8p/sorties/is3.json
% (note8p/calculs/is3_reforme.py) et note8p/sorties/is4.json (note8p/calculs/is4_reforme.py) ; figure g7 :
% note8p/calculs/is3_figure_g7.py.
\subsection{Impôt sur les sociétés : données, méthode et trajectoire}\label{ann:is}

\paragraph{Droit de référence (2026).} Taux normal de 25~\%, maintenu ; contribution sociale de 3,3~\% au-delà de
763\,000~€ d'impôt. Taux réduit de 15~\% sur les 42\,500 premiers euros de bénéfice pour les sociétés dont le chiffre
d'affaires est d'au plus 10~M€ et dont le capital, entièrement libéré, est détenu à 75~\% au moins par des personnes
physiques (CGI, art.~219~I~b ; BOI-IS-LIQ-20-10 ; service-public.fr, fiche vérifiée le 17/02/2026). Contribution
exceptionnelle sur les grandes entreprises due au titre des deux premiers exercices clos à compter du 31/12/2025, puis
éteinte (LFI 2026, art.~12 ; BOI-IS-AUT-60) : droit en vigueur, non comptée. Amortissement dégressif : coefficient 2,25
au-delà de six ans (CGI, art.~39~A). Déficits : report en avant imputable à hauteur de 1~M€ plus 50~\% du bénéfice
au-delà (BOI-IS-DEF-10-30) ; report en arrière d'un exercice dans la limite de 1~M€ (BOI-IS-DEF-20-10). Charges
financières nettes déductibles jusqu'au plus élevé de 3~M€ et de 30~\% de l'EBITDA fiscal (art.~212~bis). CIR (art.~244
quater~B) : 30~\% des dépenses jusqu'à 100~M€, 5~\% au-delà ; depuis la LF 2025 (art.~55-56, d'après Légifiscal), forfait
de fonctionnement de 40~\% des dépenses de personnel, frais de brevets exclus ; CII ramené à 20~\% et prorogé jusqu'au
31/12/2027. Créance de CIR imputée sur l'IS de N+1 à N+3, reliquat restitué en N+4 ; remboursement immédiat pour les PME
au sens européen, les jeunes entreprises innovantes, les entreprises nouvelles et en difficulté (art.~199~ter~B, d'après
le Sénat, note de présentation « Remboursements et dégrèvements », PLF~2027).

\paragraph{Données et conventions.} Bénéfices imposables de l'exercice 2025 (DGFiP, \emph{Statistiques} n°~51) : 256,8~Md€
au taux normal, 30,6~Md€ au taux réduit, soit 267,2 et 31,8~Md€ en 2027 ; +3~\% par an ensuite. Un point de taux normal :
2,99~Md€ en 2027. Investissement éligible (machines, équipements, informatique des sociétés non financières, Insee
8.212) : 84,1~Md€ en 2025, 89,1~Md€ en 2027 ; stocks de fin 2025 : 541~Md€ (équipements), 400~Md€ (bâtiments).
Hypothèses : 75~\% de l'investissement éligible fait par des sociétés imposables ; plan actuel moitié linéaire, moitié
dégressif sur 7~ans. Années d'exercice ; en caisse, un an plus tard. Dépenses fiscales : Voies et moyens, tome~II,
PLF~2026 (coûts 2024 réalisés, 2025 et 2026 prévus), portées à +3~\% par an.

\paragraph{Formules.} Nouveau plan d'amortissement : $b\,e_0 + (1-b)\,d$, où $d_k$ est le plan ordinaire et $b$ la
fraction déduite l'année de l'achat. Coût de l'année $t$ : $C_t = \tau\, s \sum_{v\le t} I_v\,[\,b_v e_0 + (1-b_v) d - d\,]_{t-v}$ ;
en régime établi, $C_t = \tau\, s\, b\, I_t\,(1 - z(g))$, avec $z(r)=\sum_k d_k/(1+r)^k$ ; coût actualisé d'une année
d'achats : $\tau\, s\, b\, I\,(1-z(r))$. Coût d'usage relatif : $(1-\tau z)/(1-\tau)$. TMEI (Devereux et Griffith) :
$\tilde p = (1-\tau z)(\rho-\pi+\delta(1+\pi))/((1-\tau)(1+\pi)) - \delta$, $\mathrm{TMEI} = (\tilde p - r)/\tilde p$ ; $r = 5~\%$,
$\pi = 2~\%$, $\delta$ = 14~\% (équipements, 7~ans), 33~\% (logiciels, 3~ans), 3,1~\% (bâtiments, 25~ans) ; $\rho = 7{,}1~\%$
(fonds propres) ou $i(1-\tau)$ (emprunt) ; mixte : un tiers d'emprunt. Dépense de R\&D passée en charges avec un crédit
non imposable au taux $c$ : $\mathrm{TMEI} = -c/(1-\tau-c)$. Effets réels : $\Delta K_j/K_j = -\varepsilon_j\, s\, \Delta\ln c_j$,
$\varepsilon$ = 1 (0,5 à 1,6) pour les équipements, 0,5 (0,25 à 1, hypothèse) pour les bâtiments ; Cobb-Douglas à emploi
donné, $\Delta Y/Y = \sum_j \sigma_j \Delta K_j/K_j$, $\sigma_j = (r+\delta_j)K_j/Y$ (3,4~\% et 1,1~\%) ; investissement de
régime établi : $(\delta_j + g)\,\Delta K_j$.

\paragraph{Taux réduit.} Rendement à plein en 2027 : 3,18~Md€ statiques (10~points sur 31,8~Md€) ; 2,15~Md€ en série
prudente (coût officiel de 2019, 2,2~Md€ pour 13~points, ramené à 10~points et actualisé), retenue pour le bouclage.
Calendrier : 20~\% pour l'exercice 2027 (moitié du rendement), 25~\% dès 2028. Perte par société : $0{,}10\min(B,
42\,500)$, au plus 4\,250~€ ; 760\,000 bénéficiaires en 2019, au moins 720\,000 sociétés en équivalent plafond en 2025.

\begin{table}[!htbp]\centering\scriptsize
\caption{Trajectoire pour le bouclage (Md€ courants, années d'encaissement ; + = recette publique)}\label{tab:is8_annuel}
\setlength{\tabcolsep}{3.2pt}
\begin{tabular}{@{}lrrrrrrrr@{}}
\toprule
& Fin du & Déduction & \multicolumn{2}{c}{Niches hors CIR} & \multicolumn{2}{c}{Suppression du CIR} & \emph{Brevets} & Solde \\
\cmidrule(lr){4-5}\cmidrule(lr){6-7}
Année & taux réduit & de 50~\% & statique & prudent & statique & prudent & \emph{(hors)} & prudent \\
\midrule
2028 & 1,08 & -- & -- & -- & -- & -- & -- & 1,08 \\
2029 & 2,22 & -- & 1,77 & 1,47 & 7,66 & 4,22 & 1,33 & 7,91 \\
2030 & 2,29 & $-$3,92 & 2,30 & 1,88 & 7,89 & 4,64 & 1,38 & 4,89 \\
2031 & 2,35 & $-$7,16 & 2,46 & 2,01 & 8,13 & 4,98 & 1,42 & 2,18 \\
2032 & 2,42 & $-$5,71 & 2,63 & 2,15 & 8,37 & 6,70 & 1,46 & 5,56 \\
2033 & 2,50 & $-$4,54 & 2,81 & 2,29 & 8,62 & 6,90 & 1,50 & 7,14 \\
2034 & 2,57 & $-$3,56 & 2,99 & 2,44 & 8,88 & 7,10 & 1,55 & 8,55 \\
2035 & 2,65 & $-$2,63 & 3,18 & 2,59 & 9,15 & 7,32 & 1,59 & 9,93 \\
2036 & 2,73 & $-$1,67 & 3,39 & 2,76 & 9,42 & 7,54 & 1,64 & 11,35 \\
2037 & 2,81 & $-$1,04 & 3,49 & 2,84 & 9,70 & 7,76 & 1,69 & 12,38 \\
2038 & 2,89 & $-$0,91 & 3,59 & 2,93 & 9,99 & 8,00 & 1,74 & 12,91 \\
2039 & 2,98 & $-$0,93 & 3,70 & 3,01 & 10,29 & 8,23 & 1,79 & 13,29 \\
2040 & 3,07 & $-$0,96 & 3,81 & 3,10 & 10,60 & 8,48 & 1,85 & 13,69 \\
\bottomrule
\end{tabular}
\par\vspace{3pt}\parbox{\linewidth}{\raggedright\scriptsize Encaissement : exercice + 1, sauf CIR (profil des créances).
Taux réduit : série prudente (20~\% pour l'exercice 2027, 25~\% ensuite). Déduction : 25~\% pour l'exercice 2029, 50~\%
ensuite, coût brut. Niches hors CIR : mécénat, crédits culture, zonages, crédit famille, dès l'exercice 2028 ; prudent :
après réaction (facteurs 0,85, 0,75, 0,80, 0,90). CIR supprimé pour les dépenses de 2028 et suivantes ; statique : coût
total de l'année ; prudent : millésimes de créances 2028 et suivants seulement (le stock antérieur est imputé ou restitué
jusqu'en 2031), décote de 20~\%. Solde prudent : colonnes 2, 3, 5 et 7. Brevets : statique, non comptés. Mémoire, non
comptés (droit en vigueur) : fin de la contribution exceptionnelle, environ $-$7~Md€ par an par rapport à 2026 ;
extinction du CII. Source : \texttt{note8p/calculs/is4\_reforme.py}.}
\end{table}

\paragraph{Déduction immédiate.} Coût permanent : 0,67~Md€ de 2027 par an (0,54 à 0,81 selon le plan actuel), 0,99~Md€ en
2040 ; coût actualisé d'une année d'achats : 1,07~Md€ à 5~\% (0,73 au taux de l'État de 3,25~\%). Le coût des premières
années est une avance d'impôt, récupérée ensuite. Véhicules utilitaires (moitié du matériel de transport, hypothèse) :
coûts multipliés par 1,25. Aucune recette de croissance n'est comptée.

\paragraph{Niches : paquet chiffré.} Coûts statiques des Voies et moyens (2026). Mécénat des entreprises (210309,
1\,730~M€, 157\,399~bénéficiaires) : réduction remplacée par la déduction du don, $1{,}73\,(1-0{,}25/0{,}60) = 1{,}01$~Md€ à
dons inchangés. Crédits culture (320121, 129, 140, 135, 142, 128, 145, 148) : 0,71~Md€, moitié en 2028. Zonages (220104,
220109, 220102, 230602, 230606, 230608, 230605, 210305) : 0,66~Md€, fermés aux entrants, extinction linéaire sur huit ans
(hypothèse). Crédit famille (210308) : 0,22~Md€. Noyau hors CIR : 2,60~Md€ (2026). Rendement réalisable prudent
(hypothèses) : mécénat $\times 0{,}85$ (taux de 40~\% au-delà de 2~M€ de dons et report des excédents sur cinq ans ; la
baisse des dons accroît au contraire la recette) ; culture $\times 0{,}75$ (tournages perdus, recettes de TVA et de
cotisations en moins) ; zonages $\times 0{,}80$ (substitution) ; crédit famille $\times 0{,}90$. Hors paquet : taxe au
tonnage (régime européen, base mobile), outre-mer (à budgétiser), PTZ et HLM (logement), SIIC et SPPICAV (transparence).
La quote-part de frais et charges et le régime mère-fille éliminent la double imposition : ils ne sont pas des niches.
Doubles comptes : les dispositifs communs IR-IS sont comptés ici en entier (le bloc IR ne les reprend pas) ; les
exonérations de fiscalité locale en zone et le mécénat des ménages relèvent d'autres blocs.

\paragraph{Suppression du CIR.} Dépenses de R\&D engagées à compter de 2028 : plus de crédit, déduction ordinaire des
charges ; équipements de recherche dans la déduction de 50~\% (déjà dans l'investissement éligible). L'ancien recentrage
(plafond au niveau du groupe et fin du taux de 5~\%, 0,97~Md€) est absorbé : pas de double compte. Coût : 7,3~Md€
exécutés en 2025, plus de 7,4~Md€ projetés en 2027 (Sénat, PLF~2027) ; 8,0~Md€ en 2026 selon les Voies et moyens du
PLF~2026 (borne haute). Créance 2023 : 7,85~Md€, dont recherche 7,32 (PME 2,24, ETI 2,10, grandes entreprises 2,98) ;
33~bénéficiaires dépassant 100~M€ de dépenses au niveau de l'intégration fiscale reçoivent 25~\% de la créance recherche
(MESRE, novembre 2025). En 2024, sur 7,7~Md€ (CIR et CII), 4,7 restitués et 3,0 imputés (Sénat, d'après la DGFiP). Profil
retenu d'un millésime de créances (hypothèse construite sur ces chiffres) : 67~\% en N+1 (remboursements immédiats,
environ 36~\%, et 80~\% des imputations), 5~\% en N+2, 3~\% en N+3, 25~\% restitués en N+4. Réalisable : millésimes 2028 et
suivants seulement, $\times 0{,}8$ (emploi de recherche perdu, soutiens recréés hors de l'impôt ; hypothèse, aucune
élasticité ne la calibre). Effet sur la R\&D : un euro de CIR produit 0,9 à 1,5~euro de R\&D dans les PME (CNEPI, cité
par le Sénat ; Bozio, Cottet et Py, 2019), aucun effet établi pour les ETI et grandes entreprises (CNEPI 2021, IPP 2021).
R\&D des PME : $-$2,1 à $-$3,4~Md€ par an ; stock $-$11 à $-$19~Md€ (dépréciation de 15~\%, croissance de 3~\%) ; PIB
potentiel $-(r+\delta)\Delta S/Y$ = $-$0,07 à $-$0,12~\%, double si le rendement social est double du rendement privé
(Bloom, Schankerman et Van Reenen, 2013). Non compté. Variante : crédit conservé pour les seules PME, recette statique
5,1~Md€ (2027). TMEI de la R\&D : $-$67~\% avec le CIR à 30~\%, nul sans.

\begin{table}[!htbp]\centering\scriptsize
\caption{Taux marginaux effectifs d'un investissement neuf (\%)}\label{tab:is8_tef}
\setlength{\tabcolsep}{4pt}
\begin{tabular}{@{}lrrrrrr@{}}
\toprule
& \multicolumn{3}{c}{Équipement} & Logiciel & \multicolumn{2}{c}{Bâtiment} \\
\cmidrule(lr){2-4}\cmidrule(l){6-7}
Configuration & fonds propres & emprunt & 1/3 dette & fonds propres & fonds propres & emprunt \\
\midrule
Aujourd'hui (25~\%) & 18,0 & $-$23,8 & 7,6 & 19,5 & 22,0 & $-$19,3 \\
Taux à 22,5~\% (comparaison) & 16,0 & $-$20,8 & 6,5 & 17,5 & 19,7 & $-$16,7 \\
Déduction de 50~\% (proposition) & 9,9 & $-$37,0 & $-$1,7 & 10,8 & 22,0 & $-$19,3 \\
Déduction intégrale (long terme) & 0,0 & $-$53,4 & $-$13,1 & 0,0 & 22,0 & $-$19,3 \\
Petite société sous 42\,500~€, aujourd'hui & 10,4 & $-$12,8 & 3,8 & 11,4 & 13,0 & $-$10,0 \\
\bottomrule
\end{tabular}
\par\vspace{3pt}\parbox{\linewidth}{\raggedright\scriptsize Bâtiments hors champ de la déduction. Dépense de R\&D passée
en charges : TMEI nul ; avec le CIR à 30~\% : $-$67~\% ; au taux de 5~\% au-delà de 100~M€ : $-$7~\%. Avantage de la dette (écart
de coût du capital entre fonds propres et emprunt) : 2,06 points aujourd'hui, 1,90 avec la déduction de 50~\%, 1,74 avec la
déduction intégrale. Niveaux non comparables à ceux de l'OCDE ; seuls les écarts comptent.}
\end{table}

\paragraph{Dette.} Coût du capital d'une machine (rendement exigé 5~\%) : fonds propres 6,10, 5,55 et 5,00~\% ; emprunt
4,04, 3,65 et 3,26~\% pour une fraction déduite tout de suite de 0, 50 et 100~\%. L'écart diminue (une déduction anticipée
vaut plus actualisée à 7,1~\% qu'à $i(1-\tau) = 5{,}3$~\%), mais le coût de l'emprunt passe sous le rendement exigé : les
intérêts nominaux restent déduits sur tout le prix alors que l'État en a déjà pris $\tau b$ à sa charge. Corrections
(TMEI pour $b = 0{,}5$) : rabot général de 21~\% des charges financières nettes pour ramener le TMEI de l'emprunt à
$-$24~\%, de 14~\% pour un TMEI mixte nul ; déduction des seuls intérêts réels ($i-\pi$) : $-$19,5~\% (emprunt), +1,8~\% (mixte) ;
au-delà du plafond de l'art.~212~bis, l'intérêt marginal n'est pas déductible et le TMEI de l'emprunt égale celui des
fonds propres (9,9~\%). Le seuil de 3~M€ de charges nettes correspond à 50 à 75~M€ de dette (taux de 6 à 4~\%), celui de
1~M€ (Espagne) à 17 à 25~M€ : l'abaisser ne touche pas l'investissement marginal des PME. Recommandation : pas de
garde-fou dans la proposition centrale ; neutralité à traiter avec la déduction intégrale.

\paragraph{Pertes.} Valeur actualisée de la déduction de 50~\% pour 100 d'achat : 1,60 pour une société imposable ;
1,07, 0,41 et 0,09 si elle reste déficitaire 1, 3 ou 5~ans (déficits non indexés, imputés sans plafond) ; nulle au-delà
de la durée du plan. Part de l'investissement éligible faite par des sociétés déficitaires : non mesurée (aucune source
croisant investissement et situation fiscale) ; repère : 113~Md€ de déficits déclarés en 2022 pour 278~Md€ de bénéfices
(DGFiP n°~51) ; l'hypothèse de 75~\% d'investissement imposable porte coûts et effets.

\paragraph{Effets réels.} Coût d'usage des équipements : $-$2,1~\% avec la déduction de 50~\%, $-$4,2~\% avec la déduction
intégrale, $-$0,5~\% avec un taux de 22,5~\% (qui baisse aussi de 1,6~\% celui des bâtiments). Capital supplémentaire
(central, bornes) : 8,4 (4,2 à 13,4)~Md€ ; 17,0 (8,5 à 27,1) pour la déduction intégrale ; 4,6 (2,3 à 8,3) pour le taux à
22,5~\%. Investissement de régime établi : +1,4~Md€ par an. PIB potentiel, effet de niveau : +0,05~\% (0,03 à 0,09), soit
1,6~Md€, pour la déduction de 50~\% ; +0,11~\% pour la déduction intégrale ; +0,02~\% pour le taux à 22,5~\%. Par Md€ de coût
économique annuel (coût actualisé d'une année d'achats, ou coût permanent du taux), 7,8~Md€ de capital pour la
déduction, 0,6 pour le taux : calcul partiel, qui ignore la localisation et les bénéfices transférés ; Chodorow-Reich,
Zidar et Zwick (2024) trouvent un écart plus faible, de même sens. Petites sociétés sous le plafond : coût d'usage des
équipements $-$0,1~\%, des bâtiments +6~\% ; non agrégé, faute de connaître leur part du capital ; effet de trésorerie non
chiffré (sensibilité non estimée en France). Ces effets ne mesurent que la déduction : la fin
du taux réduit (bâtiments et trésorerie des petites sociétés) et la suppression du CIR (R\&D des PME) jouent en sens
inverse et ne sont pas chiffrées dans le +0,05~\% ; les autres niches n'ont qu'un effet de localisation ou de secteur.

\paragraph{Long terme, non retenu.} Passer de 50 à 100~\% en 2035 : 9,4~Md€ de plus la première année, 7,5 la deuxième,
puis 0,67~Md€ par an (euros de 2027) ; d'un coup en 2027 : 14,8~Md€ ; TMEI mixte de $-$13~\%, d'où un traitement des intérêts
à prévoir. Déduction pour fonds propres et remboursement des pertes : non chiffrés (fonds propres fiscaux non disponibles
dans une source vérifiée).

\paragraph{Cas types.} Machine de 500\,000~€ achetée le 1\textsuperscript{er} juillet, 7~ans, à 25~\% : déduction de la
première année de 35\,714 à 267\,857~€ (linéaire), soit 58\,036~€ de trésorerie, rendus ensuite ; avantage actualisé de
9\,544~€ (5~\%) ou 14\,155~€ (8~\%) ; dégressif : 52\,455 et 6\,461~€. Suramortissement « industrie~4.0 » (PLF~2027, art.~9,
proposé) sur la même machine : 43\,600~€ (40~\%), 87\,200~€ (80~\%), 130\,800~€ (120~\%) actualisés, contre 8\,003~€ pour la
déduction de 50~\% (plan moyen). Société de services : perte de $0{,}10\min(B, 42\,500)$, soit 3\,000~€ à 30\,000~€ de
bénéfice ; la déduction sur 10\,000~€ d'informatique annuelle lui vaut environ 90~€.

\paragraph{Variantes écartées et hypothèses retirées.} Taux à 22,5~\% : 7,5~Md€ par an pour un effet sur le capital
domestique moindre. Recentrage du CIR (plafond au niveau du groupe, fin du taux de 5~\%), remplacé par la suppression. Retirés par parcimonie : rendement « central » du taux
réduit (fuite de 16~\%), effet de trésorerie des PME, multiplicateur d'investissement de court terme, retour par les
bénéfices déclarés, sensibilités de la part imposable ; aucune décote comportementale sur les niches (la boîte à brevets
est exclue du bouclage prudent).

\paragraph{Limites.} Part imposable de l'investissement et plan actuel (non mesurés) ; élasticité des bâtiments ;
profil d'encaissement du CIR et décotes de réaction (hypothèses) ; charges financières nettes déduites non chiffrées ; tome~II des Voies et moyens du PLF~2027 non lu ; durées d'extinction des
zonages et taux du mécénat (CGI, art.~238~bis) non revérifiés sur les textes ; issue du PLF~2027 inconnue au 18/11/2026.
